% ═══════════════════════════════════════════════════════════════
% ARBEITSBLATT: Reduktion von Silberoxid
% Thema 2.5 Redoxreaktionen | Chemie Klasse 8 | 28 BE
% ═══════════════════════════════════════════════════════════════

\documentclass[10pt, a4paper]{article}

\usepackage[utf8]{inputenc}
\usepackage[T1]{fontenc}
\usepackage[ngerman]{babel}

\usepackage{helvet}
\renewcommand{\familydefault}{\sfdefault}

\usepackage{microtype}
\usepackage[margin=1.3cm, top=1.8cm, bottom=1.2cm]{geometry}
\usepackage{enumitem}
\usepackage{tabularx}
\usepackage{booktabs}
\usepackage{array}
\usepackage{multicol}
\usepackage{tikz}
\usetikzlibrary{calc, shapes.geometric, arrows.meta}

\usepackage{amsmath, amssymb}
\usepackage{siunitx}
\sisetup{locale = DE, per-mode = symbol}

\usepackage{fancyhdr}
\usepackage{lastpage}

\usepackage{xcolor}
\usepackage{tcolorbox}
\tcbuselibrary{skins}

\usepackage{qrcode}
\usepackage[hidelinks]{hyperref}

% ═══════════════════════════════════════════════════════════════
% FARBEN
% ═══════════════════════════════════════════════════════════════

\definecolor{chemie}{HTML}{2a7a4b}
\definecolor{hellgrau}{HTML}{F5F5F5}
\definecolor{mittelgrau}{HTML}{888888}
\definecolor{dunkelgrau}{HTML}{444444}
\definecolor{silber}{HTML}{C0C0C0}
\definecolor{sauerstoff}{HTML}{c62828}
\definecolor{warnung}{HTML}{b35c00}

\colorlet{fachfarbe}{chemie}

% ═══════════════════════════════════════════════════════════════
% VARIABLEN
% ═══════════════════════════════════════════════════════════════

\newcommand{\fach}{Chemie}
\newcommand{\thema}{Reduktion von Silberoxid}
\newcommand{\klasse}{8}
\newcommand{\maxpunkte}{28}

% ═══════════════════════════════════════════════════════════════
% KOPF- UND FUSSZEILE
% ═══════════════════════════════════════════════════════════════

\pagestyle{fancy}
\fancyhf{}
\renewcommand{\headrulewidth}{0.4pt}
\renewcommand{\footrulewidth}{0pt}
\lhead{\small \fach{} Klasse \klasse{} | \thema}
\rhead{\small Name: \underline{\hspace{3cm}}}
\cfoot{\small\textcolor{mittelgrau}{Seite \thepage{} von \pageref{LastPage}}}

% ═══════════════════════════════════════════════════════════════
% BOXEN
% ═══════════════════════════════════════════════════════════════

\newtcolorbox{merkkasten}[1][]{
    colback=hellgrau,
    colframe=hellgrau,
    boxrule=0pt,
    arc=0pt,
    left=6pt, right=6pt, top=5pt, bottom=4pt,
    borderline north={2pt}{0pt}{fachfarbe},
    before skip=4pt, after skip=4pt,
    title={#1},
    fonttitle=\bfseries\small,
    coltitle=fachfarbe,
    attach boxed title to top left={yshift=-2mm, xshift=4mm},
    boxed title style={colback=hellgrau, arc=0pt, boxrule=0pt}
}

\newtcolorbox{sicherheitsbox}{
    colback=warnung!10,
    colframe=warnung,
    boxrule=1.5pt,
    arc=0pt,
    left=8pt, right=8pt, top=8pt, bottom=6pt,
    before skip=6pt, after skip=6pt
}

\newtcolorbox{versuchsbox}[1][]{
    enhanced,
    colback=fachfarbe!5,
    colframe=fachfarbe,
    boxrule=1.5pt,
    arc=0pt,
    left=8pt, right=8pt, top=12pt, bottom=6pt,
    before skip=10pt, after skip=6pt,
    title={#1},
    fonttitle=\bfseries\small,
    coltitle=white,
    attach boxed title to top left={yshift=-3mm, xshift=4mm},
    boxed title style={colback=fachfarbe, arc=0pt, boxrule=0pt, left=4pt, right=4pt}
}

% ═══════════════════════════════════════════════════════════════
% BEFEHLE
% ═══════════════════════════════════════════════════════════════

\newcommand{\luecke}[1]{\underline{\hspace{#1}}}
\newcommand{\punktebox}[1]{\hfill\small\textcolor{mittelgrau}{[\_\_/#1]}}

\newcommand{\kasten}[2]{%
    \vspace{6pt}%
    \noindent\textbf{\textcolor{fachfarbe}{Kasten #1:}} #2%
    \vspace{3pt}%
    \nopagebreak[4]%
}

\newcommand{\operator}[1]{\textbf{\textcolor{fachfarbe}{#1}}}

\newlist{teil}{enumerate}{2}
\setlist[teil,1]{label=\alph*), leftmargin=1.6em, itemsep=2pt, parsep=0pt, topsep=2pt}

\newcommand{\antwortzeilen}[1]{%
    \par\vspace{3pt}%
    \foreach \n in {1,...,#1}{%
        \noindent\rule{\linewidth}{0.3pt}\vspace{5pt}%
    }%
}

\setlength{\parindent}{0pt}
\setlength{\parskip}{2pt}

% ═══════════════════════════════════════════════════════════════
% DOKUMENT
% ═══════════════════════════════════════════════════════════════

\begin{document}

% ═══════════════════════════════════════════════════════════════
% HEADER
% ═══════════════════════════════════════════════════════════════

\begin{center}
    {\large\bfseries\textcolor{fachfarbe}{\thema}}\\[0.1em]
    {\small Arbeitsblatt | \fach{} Klasse \klasse{} | Thema 2.5 Redoxreaktionen}
\end{center}
\vspace{-0.4em}
\hrule
\vspace{0.4em}

% ═══════════════════════════════════════════════════════════════
% KASTEN 1: BEGRIFF
% ═══════════════════════════════════════════════════════════════

\begin{minipage}[t]{0.62\textwidth}
\begin{merkkasten}[Kasten 1: Der Begriff „Reduktion"]
\textbf{Reduktion} kommt vom lateinischen Wort \luecke{2.5cm} \\[0.3em]
= \luecke{1.5cm} + \luecke{1.5cm} = „\luecke{2.5cm}"\\[0.5em]
Bei der Reduktion wird ein Metalloxid zum reinen \luecke{2cm} zurückgeführt.
\end{merkkasten}
\end{minipage}%
\hfill
\begin{minipage}[t]{0.34\textwidth}
\centering
\vspace{0.3em}
\begin{tikzpicture}[scale=0.9]
    % Ag2O
    \node[draw, fill=hellgrau, minimum width=2cm, minimum height=1cm] at (0,0) {Ag$_2$O};
    % Pfeil
    \draw[->, thick, fachfarbe] (1.2,0) -- (2.3,0);
    \node[above, fachfarbe, font=\scriptsize] at (1.75,0.1) {Reduktion};
    % Ag
    \node[draw, fill=silber, minimum width=1.5cm, minimum height=1cm] at (3.5,0) {Ag};
\end{tikzpicture}
\end{minipage}

\vspace{4pt}

% ═══════════════════════════════════════════════════════════════
% KASTEN 2: SICHERHEIT
% ═══════════════════════════════════════════════════════════════

\begin{sicherheitsbox}
\textbf{\textcolor{warnung}{Kasten 2: Sicherheitshinweise}} \punktebox{4}\\[0.3em]
\begin{minipage}[t]{0.65\textwidth}
\operator{Ergänze} die Sicherheitsregeln für das Experiment:
\begin{itemize}[leftmargin=1.5em, itemsep=6pt]
    \item[\textcolor{warnung}{\textbf{!}}] \luecke{3cm} tragen
    \item[\textcolor{warnung}{\textbf{!}}] Im \luecke{2cm} arbeiten (Gase entstehen!)
    \item[\textcolor{warnung}{\textbf{!}}] Heißes Reagenzglas \luecke{3cm}
    \item[\textcolor{warnung}{\textbf{!}}] \luecke{3.5cm} verwenden
\end{itemize}
\end{minipage}%
\hfill
\begin{minipage}[t]{0.3\textwidth}
\centering
\textbf{Materialien:}\\
{\footnotesize
Silberoxid (Ag$_2$O)\\
Reagenzglas\\
Bunsenbrenner\\
Glimmspan
}
\end{minipage}
\end{sicherheitsbox}

\vspace{2pt}

% ═══════════════════════════════════════════════════════════════
% VERSUCHSPROTOKOLL
% ═══════════════════════════════════════════════════════════════

\begin{versuchsbox}[Versuchsprotokoll: Zerlegung von Silberoxid]

\textbf{Durchführung:}
\begin{enumerate}[leftmargin=2em, itemsep=1pt]
    \item Gib eine Spatelspitze Silberoxid (Ag$_2$O) ins Reagenzglas.
    \item Erhitze das Reagenzglas vorsichtig mit dem Brenner.
    \item Halte einen glühenden Holzspan an die Öffnung des Reagenzglases.
\end{enumerate}

\vspace{6pt}
\hrule
\vspace{6pt}

% ═══════════════════════════════════════════════════════════════
% KASTEN 3: BEOBACHTUNG
% ═══════════════════════════════════════════════════════════════

\textbf{Kasten 3: Beobachtung} (makroskopisch) \punktebox{6}

\vspace{4pt}
\operator{Ergänze} die Lücken:

\vspace{4pt}
\begin{tabular}{@{}p{14cm}@{}}
Das \textbf{schwarze} Silberoxid wird beim Erhitzen \luecke{3cm}.\\[0.8em]
Der glühende Holzspan \luecke{3cm} beim Hineinhalten.\\[0.8em]
Dies zeigt, dass das Gas \luecke{2.5cm} entstanden ist. \\[0.3em]
{\small (Diese Nachweisreaktion heißt \luecke{3cm})}
\end{tabular}

\vspace{6pt}
\hrule
\vspace{6pt}

% ═══════════════════════════════════════════════════════════════
% KASTEN 4: DEUTUNG + SKIZZE
% ═══════════════════════════════════════════════════════════════

\textbf{Kasten 4: Deutung} (mikroskopisch) \punktebox{8}

\vspace{4pt}

\begin{minipage}[t]{0.55\textwidth}
\operator{Erkläre}, was auf Teilchenebene passiert:

\vspace{4pt}
Beim Erhitzen wird das \luecke{2.5cm} des Silberoxids aufgebrochen.\\[0.8em]

Die \luecke{2.5cm}-Ionen lösen sich und bilden \luecke{1.5cm}-Moleküle.\\[0.8em]

Das Gas \luecke{2cm} entweicht.\\[0.8em]

Zurück bleibt reines \luecke{2cm} (metallisch glänzend).
\end{minipage}%
\hfill
\begin{minipage}[t]{0.42\textwidth}
\textbf{Skizze:} Zeichne vorher → nachher

\begin{tikzpicture}
    % Vorher
    \draw[thick] (0,0) rectangle (1.8,1.5);
    \node[below, font=\footnotesize] at (0.9,-0.1) {\textbf{vorher}};
    \node[above, font=\scriptsize] at (0.9,1.5) {Ag$_2$O};
    % Nachher
    \draw[thick] (2.5,0) rectangle (4.3,1.5);
    \node[below, font=\footnotesize] at (3.4,-0.1) {\textbf{nachher}};
    % Pfeil
    \draw[->, thick] (1.9,0.75) -- (2.4,0.75);
    % O2 Pfeil nach oben
    \draw[->, thick, sauerstoff] (4.5,1.2) -- (4.5,1.8);
    \node[right, font=\scriptsize, sauerstoff] at (4.5,1.5) {O$_2$};
\end{tikzpicture}

{\scriptsize Legende: \textcolor{gray}{$\bullet$ Ag$^+$} \quad \textcolor{sauerstoff}{$\bullet$ O$^{2-}$}}
\end{minipage}

\end{versuchsbox}

% ═══════════════════════════════════════════════════════════════
% KASTEN 5: REAKTIONSGLEICHUNG
% ═══════════════════════════════════════════════════════════════

\begin{merkkasten}[Kasten 5: Reaktionsgleichung]

\textbf{Wortgleichung:} \punktebox{2}

\vspace{2pt}
\hspace{2em} \luecke{3cm} $\longrightarrow$ \luecke{2.5cm} + \luecke{2.5cm}

\vspace{8pt}
\textbf{Formelgleichung:} \punktebox{4}

\vspace{2pt}
\hspace{2em} \luecke{0.6cm} Ag$_2$O $\longrightarrow$ \luecke{0.6cm} Ag + \luecke{0.6cm} O$_2$

\vspace{4pt}
{\small\textit{Tipp: Gleiche zuerst die Sauerstoff-Atome aus (O$_2$ = 2 Atome), dann das Silber.}}

\end{merkkasten}

\vspace{4pt}

% ═══════════════════════════════════════════════════════════════
% KASTEN 6: VERGLEICH
% ═══════════════════════════════════════════════════════════════

\kasten{6}{Oxidation und Reduktion im Vergleich} \punktebox{2}

\vspace{4pt}

\begin{tabular}{|p{4cm}|p{5.5cm}|p{5.5cm}|}
\hline
& \textbf{Oxidation} & \textbf{Reduktion} \\
\hline
\textbf{Sauerstoff wird...} & \luecke{3cm} & \luecke{3cm} \\
\hline
\textbf{Beispiel} & 4 Ag + O$_2$ $\rightarrow$ 2 Ag$_2$O & 2 Ag$_2$O $\rightarrow$ 4 Ag + O$_2$ \\
\hline
\textbf{Merkhilfe} & „Sauerstoff kommt dazu" & „Sauerstoff geht weg" \\
\hline
\end{tabular}

\vspace{6pt}

% ═══════════════════════════════════════════════════════════════
% KASTEN 7: DEFINITION
% ═══════════════════════════════════════════════════════════════

\begin{merkkasten}[Kasten 7: Definition]

\vspace{2pt}
\textbf{\textcolor{fachfarbe}{Reduktion}} = \punktebox{2}

\vspace{4pt}
\antwortzeilen{2}

\vspace{2pt}
{\small\textit{(Gegenteil der Oxidation)}}

\end{merkkasten}

% ═══════════════════════════════════════════════════════════════
% QR + PUNKTE
% ═══════════════════════════════════════════════════════════════

\vfill

\begin{minipage}[b]{0.75\textwidth}
{\small\textbf{Gesamt:} \underline{\hspace{1cm}} / \maxpunkte{} BE}
\end{minipage}%
\hfill
\begin{minipage}[b]{0.2\textwidth}
\raggedleft
\href{https://devbox42.github.io/sciencesim/chemie/kl08/03-reduktion/LP-03-reduktion.html}{\qrcode[height=1.2cm]{https://devbox42.github.io/sciencesim/chemie/kl08/03-reduktion/LP-03-reduktion.html}}\\
{\tiny Lernpfad}
\end{minipage}

\end{document}
